\documentclass[11pt]{article}
\usepackage[a4paper,margin=25mm]{geometry}
\usepackage{amsmath,amssymb,amsthm}
\usepackage[hidelinks]{hyperref}
\newtheorem{theorem}{Theorem}
\newtheorem{corollary}{Corollary}
\theoremstyle{remark}
\newtheorem{remark}{Remark}
\newcommand{\fix}{\operatorname{fix}}
\title{The Derangement Proportion of a Transitive Group\\Is at Least $1/|\Omega|$:
A Disproof of Conjecture 00000002311}
\author{GitHub: gaochengzhecpu}
\date{}
\begin{document}
\maketitle
\begin{abstract}
The source asserts that the proportion of fixed-point-free elements of
a transitive group on $\Omega$ is at least $c/|\Omega|$ with
$c=\tfrac12$, and that this constant is tight, as verified by Frobenius
groups. The inequality is true, but the tightness clause is false. For
every finite transitive group of degree $n\ge2$ the proportion is at
least $1/n$ (Cameron and Cohen). Hence no transitive group, Frobenius
or not, attains or approaches $\tfrac12\cdot\tfrac1n$, and the tight
constant is $1$; it is attained by the Frobenius group $S_3$ on three
points. The general bound, the count of derangements in every Frobenius
group, the example and the negation of the source's assertion are
verified in Lean, using Burnside's lemma from Mathlib.
\end{abstract}

\paragraph{Reading of the source.}
Let a finite group $G$ act transitively on a finite set $\Omega$ with
$n=|\Omega|$. A \emph{derangement} is an element of $G$ without fixed
points; $D$ is the set of derangements and $\delta(G)=|D|/|G|$ is the
derangement proportion. Both language versions of the source make two
assertions:
\begin{enumerate}
\item[(i)] $\delta(G)\ge c/|\Omega|$ with $c=\tfrac12$;
\item[(ii)] this constant is tight, and Frobenius groups verify the
 tightness.
\end{enumerate}
The source names only one constant, $c=\tfrac12$, and we read ``the
corrected constant'' in (ii) as this constant. A constant in a lower
bound is called tight when it cannot be replaced by a larger one. We
consider three precise forms of (ii):
\begin{enumerate}
\item[(T1)] no constant $c'>\tfrac12$ satisfies (i) for all transitive
 groups;
\item[(T2)] some Frobenius group has $\delta(G)=\tfrac12/|\Omega|$;
\item[(T3)] for every $\varepsilon>0$ some Frobenius group has
 $\delta(G)<(\tfrac12+\varepsilon)/|\Omega|$.
\end{enumerate}
We take $n\ge2$ throughout. For $n=1$ there are no derangements, so
(i) itself fails; we do not use this degenerate case. A
\emph{Frobenius group} is a transitive permutation group of degree
$n\ge2$ in which no non-identity element fixes two points and some
non-identity element fixes a point.

\begin{theorem}[Cameron--Cohen]\label{thm:cc}
Let a finite group $G$ act transitively on a finite set $\Omega$ with
$n=|\Omega|\ge2$. Then $n\,|D|\ge|G|$, that is, $\delta(G)\ge1/n$.
\end{theorem}
\begin{proof}
Write $\fix(g)$ for the number of fixed points of $g$. By Burnside's
lemma, $\sum_{g\in G}\fix(g)$ equals $|G|$ times the number of orbits
on $\Omega$, so
\begin{equation}\label{eq:one}
 \sum_{g\in G}\fix(g)=|G| .
\end{equation}
Let $G$ act on $\Omega\times\Omega$ by $g(x,y)=(gx,gy)$. A pair is fixed
by $g$ exactly when both entries are fixed, so $g$ fixes $\fix(g)^2$
pairs. The diagonal is invariant, and as $n\ge2$ there is a pair off the
diagonal; hence there are at least two orbits, and Burnside's lemma
gives
\begin{equation}\label{eq:two}
 \sum_{g\in G}\fix(g)^2\ge2\,|G| .
\end{equation}
For every $g$ we have
\begin{equation}\label{eq:pt}
 \bigl(\fix(g)-1\bigr)\bigl(\fix(g)-n\bigr)\le n\cdot\mathbf 1_D(g).
\end{equation}
Indeed, if $g\in D$ the left side is $(-1)(-n)=n$. Otherwise
$1\le\fix(g)\le n$, and the left side is at most $0$. Summing
\eqref{eq:pt} over $G$ and using \eqref{eq:one} and \eqref{eq:two},
\[
 n\,|D|\ \ge\ \sum_{g}\fix(g)^2-(n+1)\sum_{g}\fix(g)+n\,|G|
 \ \ge\ 2|G|-(n+1)|G|+n|G|=|G| .\qedhere
\]
\end{proof}

\begin{theorem}[Frobenius groups]\label{thm:frob}
Let $G$ be a finite Frobenius group of degree $n$. Then $|D|=n-1$ and
$\delta(G)=(n-1)/|G|$.
\end{theorem}
\begin{proof}
The identity fixes $n$ points and is not a derangement. A non-identity
element fixes at most one point, so it fixes exactly one point unless it
is a derangement. By \eqref{eq:one},
$|G|=n+\bigl(|G|-1-|D|\bigr)$, hence $|D|=n-1$.
\end{proof}

\begin{corollary}\label{cor:main}
The following hold.
\begin{enumerate}
\item For every finite transitive group of degree $n\ge2$,
 $\delta(G)\ge\dfrac1n>\dfrac12\cdot\dfrac1n$. In particular \textup{(i)}
 holds, also with the constant $1$ in place of $\tfrac12$.
\item \textup{(T1)}, \textup{(T2)} and \textup{(T3)} are false. So the
 conjunction of \textup{(i)} and \textup{(ii)} is false in each of the
 three readings of tightness.
\item The constant $1$ is tight in all three senses: the Frobenius
 group $S_3$ acting on three points has $\delta=\tfrac13=1/n$.
\end{enumerate}
\end{corollary}
\begin{proof}
(1) is Theorem~\ref{thm:cc}. For (2): the constant $c'=1>\tfrac12$
satisfies (i), which refutes (T1). A Frobenius group is transitive of
degree $n\ge2$, so $\delta(G)\ge1/n$; this is larger than
$\tfrac12/n$, which refutes (T2), and with $\varepsilon=\tfrac12$ it
refutes (T3). For (3): in $S_3$ a transposition fixes exactly one point
and a $3$-cycle fixes none, so $S_3$ is a Frobenius group with
$|D|=2$ and $\delta=2/6=1/3$. Thus no $c'>1$ satisfies (i), and equality
holds for a Frobenius group.
\end{proof}

\begin{remark}[Not formalised]
For a Frobenius group with point stabiliser of order $h$ one has
$|G|=nh$ and, by Theorem~\ref{thm:frob}, $n\,\delta(G)=(n-1)/h$.
Theorem~\ref{thm:cc} gives $h\le n-1$; classically $h$ divides $n-1$,
because the stabiliser acts freely on the other $n-1$ points. So the
value $n\,\delta(G)$ ranges over numbers $(n-1)/h\ge1$ and equals $1$
exactly when $h=n-1$, as for $\mathrm{AGL}(1,q)$.
Theorem~\ref{thm:cc} is due to P.~J.~Cameron and A.~M.~Cohen
(Discrete Math.\ 106/107 (1992), 135--138), who also show that for
permutation groups equality holds exactly for sharply $2$-transitive
groups. The Boston--Shalev conjecture, now a theorem of
Fulman and Guralnick, concerns a lower bound for $\delta$ independent of
$n$ for simple groups; it is not addressed here.
\end{remark}

\section*{The earlier submission and its error}
Pull request \#190 (orionsheep) addressed the same clause (ii). It
computed that the Frobenius group $\mathrm{AGL}(1,7)$ of order $42$ has
exactly $6$ derangements, so that $n\,\delta=1$, and concluded that
Frobenius groups do not verify tightness of $\tfrac12$. It was closed
without merging on 2026-10-03. The computation was correct. The
reviewer named two errors, and both are real. First, the Lean file used
Lean core only; it counted residues with a Boolean function on natural
numbers, and its decisive theorem was
\begin{center}\small\texttt{theorem not\_half\_tight : $\neg$ (6 = 3)}.\end{center}
No group, action, derangement proportion or tightness statement
occurred. Second, one Frobenius group with $n\,\delta=1$ does not show
that no other Frobenius group does better; this needs a general fact,
which that submission mentioned only in prose and only for the groups
$C_p\rtimes C_m$. The present submission proves the general fact for
all transitive groups (Theorem~\ref{thm:cc}) and for all Frobenius
groups (Theorem~\ref{thm:frob}), with Mathlib's groups and actions, and
states the tightness readings as propositions.

\section*{Formal verification and scope}
In Lean, \texttt{derangements G $\Omega$} is the set of $g$ with
$g\cdot x\ne x$ for all $x$, for a Mathlib \texttt{MulAction} (for the
full symmetric group this is the defining formula of Mathlib's own
\texttt{derangements}), and
\texttt{derangementProportion} is the quotient of the two cardinalities
in $\mathbb Q$. Transitivity is Mathlib's
\texttt{MulAction.IsPretransitive}. Burnside's lemma is Mathlib's
\begin{center}\small
\texttt{MulAction.sum\_card\_fixedBy\_eq\_card\_orbits\_mul\_card\_group},
\end{center}
applied to $\Omega$ and to $\Omega\times\Omega$ with Mathlib's product
action. The theorem \texttt{cameron\_cohen} is Theorem~\ref{thm:cc}:
for every group $G$ and action on $\Omega$ with $G$ and $\Omega$ finite,
the action transitive and $2\le|\Omega|$,
\begin{center}\small
\texttt{Nat.card G $\le$ Nat.card $\Omega$ * Nat.card (derangements G $\Omega$)}.
\end{center}
The rational forms are
\begin{center}\small
\texttt{one\_div\_card\_le\_proportion},\qquad
\texttt{half\_div\_card\_lt\_proportion}.
\end{center}

The structure \texttt{IsFrobenius G $\Omega$} has four fields:
transitivity, $2\le|\Omega|$, that a non-identity element fixing $x$ and
$y$ forces $x=y$, and the existence of a non-identity element with a
fixed point. Theorem~\ref{thm:frob} is
\begin{center}\small
\texttt{frobenius\_card\_derangements},\qquad
\texttt{frobenius\_proportion},
\end{center}
and \texttt{frobenius\_ne\_half} states that no Frobenius group has
$\delta=\tfrac12/|\Omega|$. For $S_3$, realised as
\texttt{Equiv.Perm (Fin 3)} acting on \texttt{Fin 3}, the theorems
\texttt{s3\_isFrobenius} and \texttt{s3\_proportion} prove the Frobenius
property and $\delta=1/3$.

The proposition \texttt{LowerBound c} is (i) with the constant $c$,
quantified over all finite transitive actions of degree at least two.
The three readings of tightness are
\begin{center}\small
\texttt{CannotBeImproved c},\quad
\texttt{AttainedByFrobenius c},\quad
\texttt{ApproachedByFrobenius c}.
\end{center}
\texttt{Conjecture} is the conjunction of \texttt{LowerBound (1/2)}
with the disjunction of the three readings at $c=\tfrac12$. The final
theorem \texttt{conjecture\_false} is its negation; it is proved by
refuting each of the three readings. The theorem
\texttt{corrected\_statement} proves the bound and all three readings
of tightness for $c=1$. Finite checks for $S_3$ use kernel evaluation by
\texttt{decide}; no \texttt{native\_decide}, \texttt{sorry} or custom
axiom occurs.

Scope. This is a disproof of the source's statement as a whole, through
its tightness clause; clause (i) alone is true and is proved here in a
stronger form. The actions in the Lean theorems are not assumed
faithful, so they include all transitive permutation groups. The degree
is assumed to be at least two. The Frobenius property is defined in the
file by the permutation-group condition above, because Mathlib has no
such notion; its equivalence with the kernel-and-complement description
is not used. The remark is not formalised. If ``the corrected
constant'' were meant to be a constant other than the stated
$c=\tfrac12$, the source would not say which; Corollary~\ref{cor:main}
shows that $1$ is the only constant for which the bound and the
tightness clause both hold. A reading of (ii) that does not concern the
size of the constant is not addressed.

\section*{Source and reproduction}
The exact bilingual source is included byte-for-byte as
\texttt{SOURCE.md}. Its repository path is
\begin{center}\small\texttt{conjectures/00000002311.md}.\end{center}
The project pins Lean 4.19.0 and Mathlib commit
\begin{center}\small\texttt{c44e0c8ee63ca166450922a373c7409c5d26b00b}.\end{center}
From \texttt{lean/}, run \texttt{lake build}, then
\begin{center}\small\texttt{lake env lean -DwarningAsError=true Main.lean}.\end{center}
The included public-Git manifest locks all dependency revisions.
The \texttt{verification/} directory records the fresh-build logs,
theorem-axiom reports, artifact hashes and the self-review.
No supplementary numerical computation is required.
\end{document}
